%% kapa: exam=2 round=37 year=2026 period=2 subject=appraisal_theory subject_ko=감정평가이론 seq=3 number=3 points=20
%% source: Q-Net 감정평가사 기출문제 (cst00309) artlSeq=5268128 file=raw/1008/5268128/26년 제37회 감정평가사 2차 (2교시_감정평가이론).pdf
%% conversion: kordoc 4.24.51 -> tools/gen2.py (auto); review=proofread-gpt-terra; flags=none
\begin{longquestion}{【문제 3】}
다음 자료를 참고하여 물음에 답하시오. (20점) \par
\begin{framedtext}
1. 최근 국내 대기업인 S기업의 주가상승으로 A시 아파트가격이 상승했다는 언론보도가 계속되자, 감정평가사 甲은 S기업의 주가와 A시 B동의 대표 단지인 C아파트단지 전유면적 85 m\textsuperscript{2} 아파트(이하 C아파트라 함)의 실거래가격추이를 분석하였음 \par
2. 甲은 C아파트의 최근 3년간 거래사례를 빠짐없이 모두 반영하여, S기업의 주가를 독립변수로 하는 선형회귀분석 (linear regression analysis)을 실시하였고, 분석 결과 S기업 주가와 C아파트 실거래가격 사이에는 유의한 상관 관계(correlation)가 확인되었음 \par
3. 甲은 분석 결과를 토대로 S기업의 주가가 A시 소재 전유면적 85 m\textsuperscript{2} 아파트가격에 전반적으로 영향을 미친다고 판단하고, A시 소재 전유면적 85 m\textsuperscript{2} 아파트 감정평가 시 S기업의 주가를 활용하기로 결정하였음 \par
\end{framedtext}
\medskip\noindent 물음 1) 부동산 가치형성요인 측면에서 甲의 의사결정에 대해 비판적 의견을 쓰시오. (10점) \par
\medskip\noindent 물음 2) 회귀분석의 자료선정과 결과분석 시 유의할 사항과 관련하여 甲의 의사결정에 대해 비판적 의견을 쓰시오. (단, 회귀분석의 기본가정은 충족된 것으로 가정함) (10점) \par
\end{longquestion}
